\section{Постановка задачи}

\begin{frame}{Три парадигмы обучения}
  \begin{description}
    \item[С учителем] В обучающей выборке известны ответы.
      Предсказываем ответ для новых объектов: например, диагноз по симптомам.
    \item[Без учителя] Ответов нет. Ищем структуру данных:
      например, группируем новости по сюжетам.
    \item[С подкреплением] Агент действует в среде и учится по награде,
      которая может приходить с задержкой. Например, робот учится запрыгивать на стол.
  \end{description}
\end{frame}

\begin{frame}[c]{Что предсказывает модель?}
  \begin{numrows}
        1 & Время поездки на каршеринге\\
        2 & Кликнет ли пользователь по рекламе\\
        3 & Тема научной статьи: математика, биология, психология\ldots\\
        4 & Теги ресторана: «кофе», «веранда», «живая музыка»\\
        5 & Порядок документов в поисковой выдаче\\
      \end{numrows}\par
  \medskip
  \textbf{Чем различаются ответы? Какие группы задач вы видите?}
\end{frame}

\begin{frame}[c]{Обучение с учителем: виды задач}
  \small
  \zebra{2}\small
  \begin{tabular}{>{\raggedright\arraybackslash}p{0.286\textwidth}>{\raggedright\arraybackslash}p{0.593\textwidth}}
    \toprule
    \textbf{Задача} & \textbf{Что предсказываем}\\
    \midrule
    \textbf{Регрессия} & Целевая величина --- число или вектор. Время поездки на каршеринге, спрос на товар.\\
    \textbf{Бинарная классификация} & Да или нет. Кликнет ли по рекламе, вернёт ли кредит.\\
    \textbf{Многоклассовая классификация} & Один класс из нескольких. Тема научной статьи.\\
    \textbf{Многометочная классификация} & Несколько классов сразу, они могут пересекаться. Теги ресторана.\\
    \textbf{Ранжирование} & Упорядочить объекты. Поисковая выдача по релевантности.\\
    \textbf{Структурированные ответы} & Сегментация изображения, перевод текста, генерация новых объектов.\\
    \bottomrule
  \end{tabular}
\end{frame}

\begin{frame}[c]{Определите тип задачи}
  \begin{numrows}
        1 & Предсказать курс евро к доллару на следующий день\\
        2 & Переписать разговорный текст в официальном стиле\\
        3 & Определить, котик или собачка на картинке\\
        4 & Научить робокота запрыгивать на стол из любой позы\\
        5 & Найти наборы товаров, которые покупают вместе\\
      \end{numrows}\par
\end{frame}

\begin{frame}[c]{Ответы: тип задачи}
  \small
  \zebra{2}
  \begin{tabular}{>{\raggedright\arraybackslash}p{0.267\textwidth}>{\raggedright\arraybackslash}p{0.611\textwidth}}
    \toprule
    \textbf{Задача} & \textbf{Тип}\\
    \midrule
    1. Курс евро & Регрессия: предсказываем число\\
    2. Стилизация текста & Генерация новых объектов на основе существующих\\
    3. Котик или собачка & Бинарная классификация: два класса\\
    4. Робокот & Обучение с подкреплением: агент учится по награде за действия\\
    5. Наборы товаров & Обучение без учителя: ответов нет\\
    \bottomrule
  \end{tabular}
\end{frame}
